\section{Formal status}
\label{supp:status}

\subsection{The library}

The Lean~4 library has thirteen modules for the entries G1--G13 and seven for the instances. It imports no external library, contains no \texttt{sorry} and declares no axioms of its
own.

\subsection{Axioms}

\Cref{tab:axioms} gives, for every theorem named in \Cref{tab:concordance}, the axioms reported by
\texttt{\#print axioms}. The records for the law modules are dated 2 September 2026. The records for the instance modules are dated 26 September 2026. Only the three
standard axioms of Lean's kernel occur.

{\scriptsize
\begin{longtable}{@{}>{\raggedright\arraybackslash}p{8.6cm}>{\raggedright\arraybackslash}p{3.4cm}l@{}}
\caption{Axioms of the theorems in the concordance.}\label{tab:axioms}\\
\toprule
\textbf{Declaration} & \textbf{Axioms} & \textbf{Record} \\
\midrule
\endfirsthead
\toprule
\textbf{Declaration} & \textbf{Axioms} & \textbf{Record} \\
\midrule
\endhead
\bottomrule
\endfoot
\path{Instances.Collatz.accelerated_step} & propext, Quot.sound & instances report \\
\path{Instances.Collatz.affine_ledger} & propext, Quot.sound & instances report \\
\path{Instances.Collatz.accelerated_affine_ledger} & propext, Quot.sound & instances report \\
\path{Instances.Collatz.c_one_or_two} & propext, Quot.sound & instances report \\
\path{Instances.Collatz.first_formula} & propext, Quot.sound & instances report \\
\path{Instances.Collatz.split_pair} & propext, Quot.sound & instances report \\
\path{Instances.Collatz.split_pair_next_exponent} & propext, Quot.sound & instances report \\
\path{Instances.Collatz.split_pair_prior_exponent} & propext, Quot.sound & instances report \\
\path{Instances.BinaryReverseAdd.transition_01} & propext, Quot.sound & instances report \\
\path{Instances.BinaryReverseAdd.transition_12} & propext, Quot.sound & instances report \\
\path{Instances.BinaryReverseAdd.transition_23} & propext, Quot.sound & instances report \\
\path{Instances.BinaryReverseAdd.transition_30} & propext, Quot.sound & instances report \\
\path{Instances.BinaryReverseAdd.word_transition_01} & propext, Classical.choice, Quot.sound & instances report \\
\path{Instances.BinaryReverseAdd.word_transition_12} & propext, Classical.choice, Quot.sound & instances report \\
\path{Instances.BinaryReverseAdd.word_transition_23} & propext, Classical.choice, Quot.sound & instances report \\
\path{Instances.BinaryReverseAdd.word_transition_30} & propext, Classical.choice, Quot.sound & instances report \\
\path{Instances.BinaryReverseAdd.p0_not_palindrome} & propext, Quot.sound & instances report \\
\path{Instances.BinaryReverseAdd.p1_not_palindrome} & propext & instances report \\
\path{Instances.BinaryReverseAdd.p2_not_palindrome} & propext & instances report \\
\path{Instances.BinaryReverseAdd.p3_not_palindrome} & propext & instances report \\
\path{Instances.BinaryReverseAdd.family_closed} & propext, Classical.choice, Quot.sound & instances report \\
\path{Instances.BinaryReverseAdd.family_target_free} & propext, Classical.choice, Quot.sound & instances report \\
\path{Instances.BinaryReverseAdd.pre_entry_word_transitions} & propext, Classical.choice, Quot.sound & instances report \\
\path{Instances.BinaryReverseAdd.pre_entry_orbit} & propext, Classical.choice, Quot.sound & instances report \\
\path{Instances.BinaryReverseAdd.confined_after_entry} & propext, Classical.choice, Quot.sound & instances report \\
\path{Instances.BinaryReverseAdd.orbit_22_never_palindrome} & propext, Classical.choice, Quot.sound & instances report \\
\path{Instances.Ducci.iter_pow_two} & propext, Quot.sound & instances report \\
\path{Instances.Ducci.ducci_nilpotent} & propext, Quot.sound & instances report \\
\path{Instances.Ducci.ducci_vector_nilpotent} & propext, Quot.sound & instances report \\
\path{Instances.IntegerDucci4.maximum_nonincreasing} & propext, Quot.sound & 2026-09-26 transcript \\
\path{Instances.IntegerDucci4.four_steps_even} & propext, Quot.sound & 2026-09-26 transcript \\
\path{Instances.IntegerDucci4.repeated_four_divisibility} & propext, Quot.sound & 2026-09-26 transcript \\
\path{Instances.IntegerDucci4.reaches_zero} & propext, Quot.sound & 2026-09-26 transcript \\
\path{Laws.G8OdometerAbelianization.partA_order_independent_and_odometer} & propext, Quot.sound & 2026-09-02 receipt \\
\path{Laws.G8OdometerAbelianization.order_independent_stabilization} & propext, Quot.sound & 2026-09-02 receipt \\
\path{Laws.G8OdometerAbelianization.odometer_invariance} & propext, Quot.sound & 2026-09-02 receipt \\
\path{Laws.G8OdometerAbelianization.caseB_not_abelian} & propext & 2026-09-02 receipt \\
\path{Laws.G8OdometerAbelianization.caseB_confluent_runs_share_final} & propext & 2026-09-02 receipt \\
\path{Laws.G8OdometerAbelianization.caseB_route_dependent_counter} & none & 2026-09-02 receipt \\
\path{Laws.G8OdometerAbelianization.caseD_final_states_distinct} & none & 2026-09-02 receipt \\
\path{Laws.G8OdometerAbelianization.caseD_nonconfluent_witness} & propext & 2026-09-02 receipt \\
\path{Laws.G8OdometerAbelianization.least_action} & propext, Quot.sound & 2026-09-02 receipt \\
\path{Laws.G1HiddenAmortizedSolvency.part_a_reduction} & propext, Classical.choice, Quot.sound & 2026-09-02 receipt \\
\path{Laws.G1HiddenAmortizedSolvency.part_b_discharge} & propext, Classical.choice, Quot.sound & 2026-09-02 receipt \\
\path{Laws.G2TransfiniteEscrow.no_infinite_nonterminal_run} & none & 2026-09-02 receipt \\
\path{Laws.G2TransfiniteEscrow.not_nonterminal_of_stage_coherence} & none & 2026-09-02 receipt \\
\path{Instances.G2NativeCoherence.no_infinite_nonterminal_run_native} & none & instances report \\
\path{Laws.G3AmortizedCurrency.telescoping_identity} & propext, Quot.sound & 2026-09-02 receipt \\
\path{Laws.G3AmortizedCurrency.telescoping_bound} & propext, Quot.sound & 2026-09-02 receipt \\
\path{Laws.G3AmortizedCurrency.amortized_sum_le_scale} & propext, Quot.sound & 2026-09-02 receipt \\
\path{Laws.G3AmortizedCurrency.uniform_actual_cost_bound} & propext, Quot.sound & 2026-09-02 receipt \\
\path{Laws.G4MovingCoverExhaustion.positive_schema} & none & 2026-09-02 receipt \\
\path{Laws.G4MovingCoverExhaustion.conditional_biconditional} & none & 2026-09-02 receipt \\
\path{Laws.G4MovingCoverExhaustion.fixed_package_no_go} & none & 2026-09-02 receipt \\
\path{Instances.BinaryCounter.cost_is_bit_changes} & propext & 2026-09-26 transcript \\
\path{Instances.BinaryCounter.potential_nonnegative} & propext, Quot.sound & 2026-09-26 transcript \\
\path{Instances.BinaryCounter.amortized_exact} & propext, Quot.sound & 2026-09-26 transcript \\
\path{Instances.BinaryCounter.n_increments_from_zero} & propext, Quot.sound & 2026-09-26 transcript \\
\path{Instances.TriangleSandpile.h1} & propext, Quot.sound & 2026-09-26 transcript \\
\path{Instances.TriangleSandpile.chips_drop_one} & propext, Quot.sound & 2026-09-26 transcript \\
\path{Instances.TriangleSandpile.h2} & propext, Quot.sound & 2026-09-26 transcript \\
\path{Instances.TriangleSandpile.stabilizes} & propext, Classical.choice, Quot.sound & 2026-09-26 transcript \\
\path{Instances.TriangleSandpile.canonical_stabilization} & propext, Classical.choice, Quot.sound & 2026-09-26 transcript \\
\path{Instances.MoserSpindle.no_fin3_coloring} & propext, Quot.sound & 2026-09-26 transcript \\
\path{Instances.MoserSpindle.spindle_witness} & propext, Quot.sound & 2026-09-26 transcript \\
\path{Instances.MoserSpindle.no_coloring_of_edge_image} & propext, Quot.sound & 2026-09-26 transcript \\
\path{Instances.MovingCover.residual_decreases} & propext, Quot.sound & 2026-09-26 transcript \\
\path{Instances.MovingCover.exhausted} & propext, Quot.sound & 2026-09-26 transcript \\
\path{Instances.MovingCover.finite_leak} & propext, Quot.sound & 2026-09-26 transcript \\
\path{Instances.MovingCover.no_finite_package} & propext, Quot.sound & 2026-09-26 transcript \\
\path{Laws.G5CarryHorizonConfinement.no_fixed_depth_future_sufficient} & none & 2026-09-02 receipt \\
\path{Laws.G5CarryHorizonConfinement.no_future_sufficient_depth} & none & 2026-09-02 receipt \\
\path{Laws.G5CarryHorizonConfinement.closed_pattern_persists} & none & 2026-09-02 receipt \\
\path{Laws.G5CarryHorizonConfinement.closed_pattern_confinement} & propext & 2026-09-02 receipt \\
\path{Laws.G6EndogenousNeedleGeneration.covering_schema} & none & 2026-09-02 receipt \\
\path{Laws.G6EndogenousNeedleGeneration.hole_forming_schema} & none & 2026-09-02 receipt \\
\path{Laws.G7AdversarialMobilityConfinement.escape_never_stuck} & none & 2026-09-02 receipt \\
\path{Laws.G7AdversarialMobilityConfinement.confinement_forces_trap} & propext, Classical.choice, Quot.sound & 2026-09-02 receipt \\
\path{Laws.G7AdversarialMobilityConfinement.escape_confinement_mutually_exclusive} & propext & 2026-09-02 receipt \\
\path{Laws.G9DefectEvacuationTransport.census_tail_bound} & propext & 2026-09-02 receipt \\
\path{Laws.G9DefectEvacuationTransport.census_excludes_unbounded_creation} & propext & 2026-09-02 receipt \\
\path{Laws.G9DefectEvacuationTransport.certified_transport_regime_of_assigned_nonzero} & none & 2026-09-02 receipt \\
\path{Instances.Rule184.free_flow_shift} & propext, Quot.sound & 2026-09-26 transcript \\
\path{Instances.Rule184.free_flow_invariant} & propext, Quot.sound & 2026-09-26 transcript \\
\path{Instances.Rule184.free_flow_tail} & propext, Quot.sound & 2026-09-26 transcript \\
\path{Instances.Rule184.step_jam} & propext, Quot.sound & 2026-09-26 transcript \\
\path{Instances.Rule184.step_freeA} & propext, Quot.sound & 2026-09-26 transcript \\
\path{Instances.Rule184.step_freeB} & propext, Quot.sound & 2026-09-26 transcript \\
\path{Instances.Rule184.jam_has_one_blocked} & propext & 2026-09-26 transcript \\
\path{Instances.Rule184.orbit4_blocked_zero} & propext, Quot.sound & 2026-09-26 transcript \\
\path{Instances.Rule184.plateau_low_density} & propext & 2026-09-26 transcript \\
\path{Instances.Rule184.plateau_has_jam} & propext & 2026-09-26 transcript \\
\path{Instances.Rule184.plateau_jam_count_stalls} & propext & 2026-09-26 transcript \\
\path{Instances.Rule184.low_density_four_one_step} & propext, Classical.choice, Quot.sound & 2026-09-26 transcript \\
\path{Instances.Rule184Conservation.cars_conserved} & propext, Quot.sound & 2026-09-26 transcript \\
\path{Instances.Rule184Small.clears_up_to_eight} & propext, Classical.choice, Quot.sound & 2026-09-26 transcript \\
\path{Laws.G10LossfulPersistence.boundary_persistence} & propext & 2026-09-02 receipt \\
\path{Laws.G10LossfulPersistence.checksum_lossful} & none & 2026-09-02 receipt \\
\path{Laws.G10LossfulPersistence.checksum_protect} & propext & 2026-09-02 receipt \\
\path{Laws.G10LossfulPersistence.checksum_regen} & propext & 2026-09-02 receipt \\
\path{Laws.G10LossfulPersistence.checksum_boundary_persistence} & propext & 2026-09-02 receipt \\
\path{Laws.G11GlobalAntiSymmetry.hierarchy_forces_aperiodic} & none & 2026-09-02 receipt \\
\path{Laws.G11GlobalAntiSymmetry.g11_global_anti_symmetry} & none & 2026-09-02 receipt \\
\path{Laws.G12FiniteWitnessRadiation.no_global_k_coloring} & none & 2026-09-02 receipt \\
\path{Laws.G12FiniteWitnessRadiation.orbit_radiation} & none & 2026-09-02 receipt \\
\path{Laws.G12FiniteWitnessRadiation.finite_witness_iff_no_global_coloring} & none & 2026-09-02 receipt \\
\path{Laws.G13ThinOrbitSaturation.relative_density_theorem} & propext & 2026-09-02 receipt \\
\path{Instances.ThinOrbit.power_members_iff} & propext, Quot.sound & 2026-09-26 transcript \\
\path{Instances.ThinOrbit.power_count_sublinear} & propext, Classical.choice, Quot.sound & 2026-09-26 transcript \\
\path{Instances.ThinOrbit.powers_infinite} & propext, Quot.sound & 2026-09-26 transcript \\
\path{Instances.ThinOrbit.powers_relative_zero} & propext, Classical.choice, Quot.sound & 2026-09-26 transcript \\
\path{Instances.ThinOrbit.counted_exceptions} & none & 2026-09-26 transcript \\
\path{Instances.ThinOrbit.concrete_zero_density_failure} & propext, Classical.choice, Quot.sound & 2026-09-26 transcript \\

\end{longtable}
}

\subsection{What each Lean statement covers}

\Cref{tab:clauses} records, for each law, which parts of the printed statement the Lean declarations
establish and which are argued on paper or imported.

{\small
\begin{longtable}{@{}l>{\raggedright\arraybackslash}p{6.4cm}>{\raggedright\arraybackslash}p{7.0cm}@{}}
\caption{Printed statements and their Lean coverage.}\label{tab:clauses}\\
\toprule
\textbf{Law} & \textbf{In Lean} & \textbf{On paper or imported} \\
\midrule
\endfirsthead
\toprule
\textbf{Law} & \textbf{In Lean} & \textbf{On paper or imported} \\
\midrule
\endhead
\bottomrule
\endfoot
G1 & the reformulation; ghost discharge with both hypotheses; the Collatz ledger and split pairs,
with the map defined on all natural numbers & the $2$-adic completion and the ghost; oddness of the
split pairs; the $2^L - 1$ family; the cycle condition \\
G2 & no run is nonterminal at every step, in the native form comparing $T_k(x_k)$ with $x_k$ & the
printed conclusion that some $x_k$ lies in $A$ (classical extraction); Goodstein sequences; Kirby--Paris \\
G3 & identity and bounds over integer costs; the binary counter & rational costs; the doubling
array \\
G4 & the three statements of \Cref{thm:G4}; exhaustion and the finite leak for initial intervals & greedy unit fractions; Erd\H{o}s--Straus results \\
G5 & the no-go template; confinement; the four-phase family and the orbit of $22$ & the
unbounded-horizon hypothesis for base $2$ \\
G6 & covering from ``a positive hole count later decreases''; hole formation from one-step
propagation of non-visitation & the bridge from rate inequalities (open); windows and blocks; the
designed variants \\
G7 & the three parts, confinement in rank form; exclusion given a state in the escape invariant &
other confinement arguments; the angel thresholds \\
G8 & order independence of stable legal runs from a start with no infinite legal run; the two
examples; the narrower least-action statement; the three-site sandpile & the least-action principle;
sandpiles on arbitrary connected graphs \\
G9 & the census bound for a sequence of natural numbers, in the form that the count is not
unbounded on every tail; packaging of a supplied transporter record
with one assigned defect; the rule-184 partial results of \Cref{prop:rule184}, including car conservation and evacuation
for $n \le 8$ & the evacuation template beyond packaging; evacuation for every low-density ring (open); the defect ledger; the absorption
hypothesis; the recurrence equality; rule~184 \\
G10 & the persistence theorem; the checksum cascade; Ducci nilpotency over $\mathbb{F}_2$; the
integer Ducci collapse for $n = 4$ & the integer collapse for other powers of two; Gilbreath \\
G11 & aperiodicity from the hierarchy and forced breaks, with admissibility and ``breaks $p$'' as
arbitrary predicates & finiteness of the regions $R_p$; the tile
sets \\
G12 & restriction; orbit radiation; the equivalence for one $k$ under the compactness hypothesis &
the compactness theorem; the plane embedding of the Moser spindle, whose colouring obstruction is
checked; the Hadwiger--Nelson witnesses \\
G13 & the relative-density theorem, conditional on its two count hypotheses, with $o(\cdot)$
written as ``for every $j$, eventually $j$ times the count is at most'' and $c$ rational; the powers of two as
an infinite missed family of density zero & the exponent bound; the application to a thick target;
the reciprocity obstruction; all saturation theorems \\
Bridge & the underlying theorems for parts 1 and 2 and for the budgeted and eventually-zero cases
of part 3 & the status readings; the summable and contractive cases \\
\end{longtable}
}

\subsection{Hypotheses carried but not used}

\Cref{tab:unused} lists the hypotheses that appear in a Lean statement and are not used by its
proof.

{\small
\begin{longtable}{@{}>{\raggedright\arraybackslash}p{7.6cm}>{\raggedright\arraybackslash}p{6.4cm}@{}}
\caption{Unused hypotheses.}\label{tab:unused}\\
\toprule
\textbf{Declaration} & \textbf{Unused hypothesis} \\
\midrule
\endfirsthead
\toprule
\textbf{Declaration} & \textbf{Unused hypothesis} \\
\midrule
\endhead
\bottomrule
\endfoot
\path{Laws.G10LossfulPersistence.boundary_persistence} & the map is not injective \\
\path{Laws.G11GlobalAntiSymmetry.hierarchy_forces_aperiodic} & the zero period acts as the
identity \\
\path{Laws.G3AmortizedCurrency.telescoping_bound}, \path{Laws.G3AmortizedCurrency.uniform_actual_cost_bound} & actual costs are nonnegative \\
\path{Laws.G8OdometerAbelianization.least_action} & abelian compatibility; stability of the final
configuration of the run whose counters are bounded \\
\end{longtable}
}
